\documentclass[11pt]{article}
\usepackage[a4paper,margin=25mm]{geometry}
\usepackage[T1]{fontenc}
\usepackage[utf8]{inputenc}
\usepackage{amsmath,amssymb,amsthm}
\usepackage[hidelinks]{hyperref}
\newtheorem{proposition}{Proposition}
\begin{document}
\title{Sequences and limits\\\large A LeanMirror companion}
\author{A companion adapted from Jiří Lebl}
\date{}
\maketitle
Adapted from Basic Analysis I. Selected definitions and the two propositions on uniqueness and boundedness. Other examples, exercises, figures and subsequent sections are not included or certified. Distributed under CC-BY-SA-4.0.

Original source: \url{https://github.com/jirilebl/ra}.\par
Revision: \texttt{e21ec524ca7d54f800c693b948020c188d21d01f}.
\section*{Definitions and notation}
\hypertarget{def-indices}{}\paragraph{Index convention}
Natural numbers start at 1, as in the source. A sequence is a function on these positive integers.
\[ \mathbb{N}=\{1,2,3,\ldots\},\qquad x:\mathbb{N}\to\mathbb{R}. \]
\hypertarget{def-sequence}{}\paragraph{Sequence and boundedness}
A sequence is bounded if there exists a real number B such that the absolute value of every term is at most B.
there exists $B\in \mathbb{R}$ such that for every $n\in \mathbb{N}$, $\left|\hyperlink{def-sequence}{x_{n}}\right| \le B$.
\hypertarget{def-convergence}{}\paragraph{Convergence and limit}
A sequence converges to x if, for every epsilon > 0, there exists an M such that all terms with n at least M are within epsilon of x. A sequence is convergent when it has a limit.
for every $\epsilon\in \mathbb{R}$ satisfying $0 < \epsilon$, there exists $M\in \mathbb{N}$ such that for every $n\in \mathbb{N}$ satisfying $\hyperlink{def-indices}{n\ge M}$, $\left|\hyperlink{def-sequence}{x_{n}} - x\right| < \epsilon$.
\hypertarget{def-maximum}{}\paragraph{Finite maximum}
Take the maximum over the terms with n < M together with 1 + |x|. The initial set is empty when M = 1; the additional element keeps the maximum defined.
\[ B=\max\left(\{|x_n|:n\in\mathbb{N},\ n<M\}\cup\{1+|x|\}\right). \]
\begin{proposition}A convergent sequence has a unique limit.

For every $\{x_n\}_{n=1}^{\infty}\in \mathbb{R}^{\mathbb{N}}$, for every $x\in \mathbb{R}$, for every $y\in \mathbb{R}$ satisfying $\left(\hyperlink{def-convergence}{\lim_{n\to\infty} x_n = x} \;\land\; \hyperlink{def-convergence}{\lim_{n\to\infty} x_n = y}\right)$, $y = x$.\end{proposition}
\begin{proof}Let $\{x_n\}_{n=1}^{\infty}$.
Let $x\in \mathbb{R}$.
Let $y\in \mathbb{R}$.
Suppose $\left(\hyperlink{def-convergence}{\lim_{n\to\infty} x_n = x} \;\land\; \hyperlink{def-convergence}{\lim_{n\to\infty} x_n = y}\right)$.
It suffices to show that for every $\epsilon\in \mathbb{R}$ satisfying $0 < \epsilon$, $\left|y - x\right| < \epsilon$.
Let $\epsilon\in \mathbb{R}$.
Suppose $0 < \epsilon$.
Choose $M_1\in \mathbb{N}$ such that for every $n\in \mathbb{N}$ satisfying $\hyperlink{def-indices}{n\ge M_1}$, $\left|\hyperlink{def-sequence}{x_{n}} - x\right| < \frac{\epsilon}{2}$.
Choose $M_2\in \mathbb{N}$ such that for every $n\in \mathbb{N}$ satisfying $\hyperlink{def-indices}{n\ge M_2}$, $\left|\hyperlink{def-sequence}{x_{n}} - y\right| < \frac{\epsilon}{2}$.
Choose $n\in \mathbb{N}$ such that $\left(\hyperlink{def-indices}{n\ge M_1} \;\land\; \hyperlink{def-indices}{n\ge M_2}\right)$.
Estimate:
\[ \begin{aligned}\left|y - x\right| &= \left|\left(\hyperlink{def-sequence}{x_{n}} - x\right) - \left(\hyperlink{def-sequence}{x_{n}} - y\right)\right| \\ &\le \left|\hyperlink{def-sequence}{x_{n}} - x\right| + \left|\hyperlink{def-sequence}{x_{n}} - y\right| \\ &< \frac{\epsilon}{2} + \frac{\epsilon}{2} \\ &= \epsilon\end{aligned}. \]
Hence $\left|y - x\right| < \epsilon$.
Consequently, $y = x$.
\end{proof}
\begin{proposition}A convergent sequence is bounded.

For every $\{x_n\}_{n=1}^{\infty}\in \mathbb{R}^{\mathbb{N}}$, for every $x\in \mathbb{R}$ satisfying $\hyperlink{def-convergence}{\lim_{n\to\infty} x_n = x}$, $\hyperlink{def-sequence}{\{x_n\}_{n=1}^{\infty}\text{ is bounded}}$.\end{proposition}
\begin{proof}Let $\{x_n\}_{n=1}^{\infty}$.
Let $x\in \mathbb{R}$.
Suppose $\hyperlink{def-convergence}{\lim_{n\to\infty} x_n = x}$.
Choose $M\in \mathbb{N}$ such that for every $n\in \mathbb{N}$ satisfying $\hyperlink{def-indices}{n\ge M}$, $\left|\hyperlink{def-sequence}{x_{n}} - x\right| < 1$.
Estimate, for every $n\in \mathbb{N}$ satisfying $\hyperlink{def-indices}{n\ge M}$:
\[ \begin{aligned}\left|\hyperlink{def-sequence}{x_{n}}\right| &= \left|\left(\left(\hyperlink{def-sequence}{x_{n}} - x\right) + x\right)\right| \\ &\le \left|\hyperlink{def-sequence}{x_{n}} - x\right| + \left|x\right| \\ &< 1 + \left|x\right|\end{aligned}. \]
Choose $B\in \mathbb{R}$ such that $B = \hyperlink{def-maximum}{\max\{ |x_1|,\ldots,|x_{M-1}|,1+|x|\}}$.
Then for every $n\in \mathbb{N}$, $\left|\hyperlink{def-sequence}{x_{n}}\right| \le B$.
Hence $\hyperlink{def-sequence}{\{x_n\}_{n=1}^{\infty}\text{ is bounded}}$.
\end{proof}
\section*{Verification and scope}
The formal interpretation of every proof node is checked by Lean. The argument structure and all displayed formulas are generated from the checked expressions. The original wording can be compared in the interactive companion. Rendering, vocabulary and editorial alignment remain inspectable software and human review boundaries.
\end{document}
